\documentclass[handout]{beamer}
\mode<presentation>
\usepackage{xcolor}
\usepackage[toc]{multitoc}
\usepackage{color}
\usepackage{graphicx}
\graphicspath{{../buku_ajar/buku_ajar_logika_matematika/figures/}}
\usepackage{tikz}
\usetikzlibrary{shapes,arrows,positioning,calc}
\usepackage{listings}
\usepackage{lmodern}
\usetheme[secheader]{Boadilla}
\usecolortheme{whale}
\setbeamertemplate{caption}[numbered]
\setbeamertemplate{bibliography item}[text]
\setbeamertemplate{navigation symbols}{}

\theoremstyle{definition}
\newtheorem{definisi}[theorem]{Definisi}
\theoremstyle{example}
\newtheorem{contoh}[theorem]{Contoh}

\renewcommand*{\multicolumntoc}{2}

\setbeamertemplate{footline}
{
  \leavevmode%
  \hbox{%
  \begin{beamercolorbox}[wd=.333333\paperwidth,ht=2.25ex,dp=1ex,center]{author in head/foot}%
    \usebeamerfont{author in head/foot}\insertshortauthor%
  \end{beamercolorbox}%
  \begin{beamercolorbox}[wd=.433333\paperwidth,ht=2.25ex,dp=1ex,center]{title in head/foot}%
    \usebeamerfont{title in head/foot}\insertshorttitle
  \end{beamercolorbox}%
  \begin{beamercolorbox}[wd=.233333\paperwidth,ht=2.25ex,dp=1ex,right]{date in head/foot}%
    \usebeamerfont{date in head/foot}\insertshortdate{}\hspace*{2em} \insertframenumber{} / \inserttotalframenumber\hspace*{2ex} 
  \end{beamercolorbox}}%
  \vskip0pt%
}

\subtitle[Logika Matematika]{Logika Matematika}
\title[Tabel Kebenaran, Tautologi, Kontradiksi]{Tabel Kebenaran, Tautologi, Kontradiksi, dan Kontingensi}
\author{Cecep Suwanda, S.Si., M.Kom.}
\date{Maret 2025}
\institute{Universitas Bale Bandung}

\begin{document}

\maketitle

\section*{Daftar Isi}
\begin{frame}{Daftar Isi}
\frametitle{Daftar Isi}
\tableofcontents
\end{frame}

% ============================================================
% SECTION 4-1: Tautologi, Kontradiksi, Kontingensi
% ============================================================
\section{Tautologi, Kontradiksi, dan Kontingensi}

\begin{frame}
\frametitle{Klasifikasi Proposisi Majemuk}
Proposisi majemuk dikelompokkan menjadi tiga jenis berdasarkan pola nilai kebenarannya:
\begin{itemize}
  \item \textbf{Tautologi} --- selalu benar
  \item \textbf{Kontradiksi} --- selalu salah
  \item \textbf{Kontingensi} --- bergantung kombinasi
\end{itemize}
\end{frame}

\begin{frame}
\frametitle{Tautologi}
\begin{definisi}[Tautologi]
Proposisi majemuk yang \textbf{selalu bernilai benar (T)} untuk setiap kombinasi nilai kebenaran proposisi atomik penyusunnya.
\end{definisi}
\begin{contoh}
$p \vee \neg p$ (hukum tiada jalan tengah) --- selalu benar, apa pun nilai $p$.
\end{contoh}
\begin{block}{Aplikasi}
Ekspresi selalu benar tidak memberikan penyaringan kondisi yang efektif (hasil evaluasi selalu lolos).
\end{block}
\end{frame}

\begin{frame}
\frametitle{Kontradiksi}
\begin{definisi}[Kontradiksi]
Proposisi majemuk yang \textbf{selalu bernilai salah (F)} untuk setiap kombinasi nilai kebenaran proposisi atomik penyusunnya.
\end{definisi}
\begin{contoh}
$p \wedge \neg p$ (hukum non-kontradiksi) --- selalu salah; $p$ dan $\neg p$ tidak mungkin benar bersamaan.
\end{contoh}
\begin{block}{Aplikasi}
Kondisi kontradiktif dapat menyebabkan \textit{dead code} (cabang program tidak pernah dieksekusi).
\end{block}
\end{frame}

\begin{frame}
\frametitle{Kontingensi}
\begin{definisi}[Kontingensi]
Proposisi majemuk yang nilai kebenarannya \textbf{bergantung} pada kombinasi nilai proposisi atomik; pada sebagian kombinasi bernilai benar, pada kombinasi lain bernilai salah.
\end{definisi}
\begin{contoh}
$p \rightarrow q$, $(p \vee q) \wedge \neg r$ --- sebagian besar ekspresi praktik termasuk kontingensi.
\end{contoh}
\end{frame}

% ============================================================
% SECTION 4-2: Ekuivalensi Logis
% ============================================================
\section{Ekuivalensi Logis Melalui Tabel Kebenaran}

\begin{frame}
\frametitle{Definisi Ekuivalen Logis}
\begin{definisi}[Ekuivalen Logis]
Dua proposisi $P$ dan $Q$ \textbf{ekuivalen logis} ($\equiv$) apabila kolom nilai kebenaran keduanya identik untuk seluruh kombinasi nilai variabel.\\
Secara aljabar: $(P \leftrightarrow Q)$ merupakan \textbf{tautologi}.
\end{definisi}
\end{frame}

\begin{frame}
\frametitle{Verifikasi Ekuivalensi: Hukum De Morgan}
\begin{block}{Pertanyaan}
Apakah $\neg(p \wedge q)$ ekuivalen dengan $\neg p \vee \neg q$?
\end{block}
\begin{center}
\small
\begin{tabular}{|c|c||c|c||c|}
\hline
$p$ & $q$ & $\neg(p \wedge q)$ & $\neg p \vee \neg q$ & \\ \hline
T & T & F & F & \\ \hline
T & F & T & T & \\ \hline
F & T & T & T & \\ \hline
F & F & T & T & \\ \hline
\end{tabular}
\end{center}
Kedua kolom identik $(F,T,T,T)$ $\Rightarrow$ $\neg(p \wedge q) \equiv \neg p \vee \neg q$
\end{frame}

% ============================================================
% SECTION 4-3: Hukum Dasar Aljabar Proposisi
% ============================================================
\section{Hukum Dasar Aljabar Proposisi}

\begin{frame}
\frametitle{Hukum Dominasi dan Identitas}
\begin{block}{Hukum Dominasi}
$p \vee \mathbf{T} \equiv \mathbf{T}$ \quad $p \wedge \mathbf{F} \equiv \mathbf{F}$
\end{block}
\begin{block}{Hukum Identitas}
$p \wedge \mathbf{T} \equiv p$ \quad $p \vee \mathbf{F} \equiv p$
\end{block}
\end{frame}

\begin{frame}
\frametitle{Hukum Negasi dan Involusi}
\begin{block}{Hukum Negasi}
$p \vee \neg p \equiv \mathbf{T}$ \quad (Tiada Jalan Tengah)\\
$p \wedge \neg p \equiv \mathbf{F}$ \quad (Non-Kontradiksi)
\end{block}
\begin{block}{Hukum Involusi}
$\neg(\neg p) \equiv p$
\end{block}
\end{frame}

\begin{frame}
\frametitle{Hukum Idempoten dan Absorpsi}
\begin{block}{Hukum Idempoten}
$p \vee p \equiv p$ \quad $p \wedge p \equiv p$
\end{block}
\begin{block}{Hukum Absorpsi}
$p \vee (p \wedge q) \equiv p$\\
$p \wedge (p \vee q) \equiv p$
\end{block}
\end{frame}

% ============================================================
% SECTION 4-4: Hukum Lanjutan
% ============================================================
\section{Hukum Lanjutan: Komutatif, Distributif, De Morgan}

\begin{frame}
\frametitle{Hukum Komutatif dan Asosiatif}
\begin{block}{Komutatif}
$p \vee q \equiv q \vee p$ \quad $p \wedge q \equiv q \wedge p$
\end{block}
\begin{block}{Asosiatif}
$(p \vee q) \vee r \equiv p \vee (q \vee r)$\\
$(p \wedge q) \wedge r \equiv p \wedge (q \wedge r)$
\end{block}
\end{frame}

\begin{frame}
\frametitle{Hukum Distributif}
\begin{align*}
p \vee (q \wedge r) &\equiv (p \vee q) \wedge (p \vee r) \\
p \wedge (q \vee r) &\equiv (p \wedge q) \vee (p \wedge r)
\end{align*}
Dapat digunakan dua arah: ekspansi atau faktorisasi.
\end{frame}

\begin{frame}
\frametitle{Hukum De Morgan}
\begin{block}{Teorema}
Negasi masuk ke dalam kurung: AND $\leftrightarrow$ OR, tiap komponen dinegasi.
\end{block}
\begin{align*}
\neg(p \wedge q) &\equiv \neg p \vee \neg q \\
\neg(p \vee q) &\equiv \neg p \wedge \neg q
\end{align*}
\end{frame}

% ============================================================
% SECTION 4-5: Transformasi Implikasi dan Biimplikasi
% ============================================================
\section{Transformasi Implikasi dan Biimplikasi}

\begin{frame}
\frametitle{Transformasi Implikasi}
\begin{block}{Ekuivalensi}
\begin{align*}
p \rightarrow q &\equiv \neg p \vee q \\
p \rightarrow q &\equiv \neg(p \wedge \neg q)
\end{align*}
\end{block}
\begin{contoh}
``Jika Anda menyumbang darah ($p$), maka saya memberi apresiasi ($q$).''\\
$\equiv$ ``Anda tidak menyumbang ($\neg p$) ATAU saya memberi apresiasi ($q$)''
\end{contoh}
\end{frame}

\begin{frame}
\frametitle{Transformasi Biimplikasi}
\begin{align*}
p \leftrightarrow q &\equiv (p \rightarrow q) \wedge (q \rightarrow p) \\
&\equiv (\neg p \vee q) \wedge (\neg q \vee p) \\
p \leftrightarrow q &\equiv (p \wedge q) \vee (\neg p \wedge \neg q)
\end{align*}
\end{frame}

\begin{frame}
\frametitle{Pola Penyederhanaan Campuran}
\begin{align*}
p \vee (\neg p \wedge q) &\equiv (p \vee q) \\
p \wedge (\neg p \vee q) &\equiv (p \wedge q)
\end{align*}
\textbf{Bukti} pertama: distributif $\to$ negasi $\to$ identitas.
\end{frame}

% ============================================================
% SECTION 4-6: Strategi Penyederhanaan
% ============================================================
\section{Strategi Penyederhanaan Bertahap}

\begin{frame}
\frametitle{Urutan Operasi Penyederhanaan}
\begin{enumerate}
  \item \textbf{Ubah implikasi/biimplikasi}: $\rightarrow$, $\leftrightarrow$ $\to$ $\wedge$, $\vee$, $\neg$
  \item \textbf{Dorong negasi ke dalam}: De Morgan, negasi ganda
  \item \textbf{Sederhanakan pola langsung}: identitas, dominasi, idempoten, absorpsi
  \item \textbf{Gunakan distributif}: ekspansi atau faktorisasi
\end{enumerate}
\end{frame}

\begin{frame}
\frametitle{Contoh Penyederhanaan}
Sederhanakan: $\neg(p \vee \neg q) \vee (p \wedge q \wedge \neg r)$
\begin{enumerate}
  \item De Morgan: $(\neg p \wedge q) \vee (p \wedge (q \wedge \neg r))$
  \item Involusi: $\neg(\neg q) \equiv q$
  \item Distributif dan negasi: $q \wedge (\neg p \vee \neg r)$
\end{enumerate}
\textbf{Hasil:} $\mathbf{q \wedge (\neg p \vee \neg r)}$
\end{frame}

% ============================================================
% SECTION 4-7: Studi Kasus
% ============================================================
\section{Studi Kasus: Optimasi Kode dan Gerbang Logika}

\begin{frame}[fragile]
\frametitle{Kode Asli (Ruwet)}
\begin{lstlisting}[basicstyle=\ttfamily\scriptsize, frame=single]
def cek_validitas_pinjaman(a_kerja, b_jaminan, c_bebas_utang):
  if not (not a_kerja and b_jaminan):
      if a_kerja and (b_jaminan or not c_bebas_utang):
          return True
  return False
\end{lstlisting}
Logika: $\neg(\neg A \wedge B) \wedge (A \wedge (B \vee \neg C))$
\end{frame}

\begin{frame}
\frametitle{Penyederhanaan Logika}
\begin{align*}
\neg(\neg A \wedge B) \wedge A \wedge (B \vee \neg C) &\xrightarrow{\text{De Morgan}} (A \vee \neg B) \wedge A \wedge (B \vee \neg C) \\
&\xrightarrow{\text{Komutatif}} A \wedge (A \vee \neg B) \wedge (B \vee \neg C) \\
&\xrightarrow{\text{Absorpsi}} \mathbf{A} \wedge (B \vee \neg C)
\end{align*}
\end{frame}

\begin{frame}[fragile]
\frametitle{Kode Refactoring}
\begin{lstlisting}[basicstyle=\ttfamily\scriptsize, frame=single]
def cek_validitas_pinjaman(a_kerja, b_jaminan, c_bebas_utang):
  return a_kerja and (b_jaminan or not c_bebas_utang)
\end{lstlisting}
Logika lebih ringkas, mudah dibaca, mudah diuji.
\end{frame}

\begin{frame}
\frametitle{Optimasi Gerbang Logika}
\begin{block}{Rancangan Awal}
$Q_{output} = (A \wedge B) \vee (A \wedge \neg B)$
\end{block}
\begin{align*}
&\equiv A \wedge (B \vee \neg B) \quad \text{(Distributif)} \\
&\equiv A \wedge \mathbf{T} \quad \text{(Negasi)} \\
&\equiv \mathbf{A} \quad \text{(Identitas)}
\end{align*}
Rangkaian cukup mengikuti sinyal $A$ saja.
\end{frame}

% ============================================================
% RANGKUMAN
% ============================================================
\section{Rangkuman}

\begin{frame}
\frametitle{Rangkuman}
\begin{itemize}
  \item Tautologi (selalu T), kontradiksi (selalu F), kontingensi (bergantung).
  \item Ekuivalensi logis: kolom identik $\Leftrightarrow$ $(P \leftrightarrow Q)$ tautologi.
  \item Hukum dasar: dominasi, identitas, negasi, involusi, idempoten, absorpsi.
  \item Hukum lanjutan: komutatif, asosiatif, distributif, De Morgan.
  \item Implikasi: $p \rightarrow q \equiv \neg p \vee q$.
  \item Prosedur 4 langkah penyederhanaan; aplikasi refactoring kode dan optimasi sirkuit.
\end{itemize}
\end{frame}

\begin{frame}
\frametitle{Kata Kunci}
\begin{center}
Tautologi $\bullet$ Kontradiksi $\bullet$ Kontingensi $\bullet$ Ekuivalensi Logis\\
Tabel Kebenaran $\bullet$ Hukum Aljabar $\bullet$ De Morgan $\bullet$ Refactoring
\end{center}
\end{frame}

\end{document}
